\section{Questions for Shashank Jayakumar regarding his studies of directionality in PROSPECT}

\begin{enumerate}
    \item Extracting raw counts from Fig 3 (Phys Rev paper), similarly how it was presented in AAP 2023 (by Bryce, slide 18) --- Table 4.2 in the thesis?
    \item How pin-wheels affect directionality; in other words, there is a slight rotation of segments?
    \item The reason behind excluding diagonal segments --- how does the main result/formalism change if they were included?
    \item What does calculated uncertainty physically mean?
    \item Is it a fair assumption --- $\Delta \theta_\nu = \Delta \theta_n$? How does intrinsic 26-degree spread play a role (spread in neutron angle after first scatter)?
    \item Can the formalism work at small event counts?
    \item At $N \to \infty$, the uncertainty goes to zero --- is it physical? Should it be approaching non-zero finite limit? (think of large cores)
    \item The rationale behind including $n_0$ in the denominator? It seems to be the main factor behind the smallness of angular uncertainty. {\it ``$n_0$ is the background-subtracted numberof IBD pairs where the prompt and delayed events
are in the same segment''}
\item 
\end{enumerate}

